\uuid{lmQk}
\titre{Écriture algébrique}
\auteur{bodin ; hanani ; mzari}
\organisation{exo7}
\licence{CC BY-NC-SA 4.0}
\datecreate{2020-01-20}
\niveau{L1}
\module{Algèbre}
\chapitre{Nombres complexes}
\sousChapitre{Forme cartésienne, forme polaire}
\difficulte{1}
\theme{Écritures algébrique et géométrique, écriture algébrique}
\exo7qcmid{lille-1.nombres-complexes.001}

\contenu{
\begin{question}
\qclassification{104.01}
\qtags{motcle=écriture algébrique}
Soit $z=(1-2i)^2$. Quelles sont les assertions vraies ?
\begin{answers}
    \bad{$z=5-4i$}

    \good{$z=-3-4i$}

    \bad{Le conjugué de $z$ est : $\overline{z}=3+4i$.}

    \good{Le module de $z$ est $5$.}
\end{answers}
\begin{explanations}
On développe $(1-2i)^2$. Si $z=a+ib, a,b \in \Rr, \overline{z}=a-ib$  et $|z|^2= a^2+b^2$.
\end{explanations}
\end{question}
}
